\documentclass[12pt,a4paper]{article}
\usepackage{ugmath}
\usepackage{float}
\usepackage{placeins}
\usepackage[labelfont=bf,labelsep=space]{caption}
\newcommand{\paren}[1]{\left( #1 \right)}
\newcommand{\alumno}{Ricardo León Martínez}
\newcommand{\materia}{Fisica I}
\newcommand{\profesor}{Lucero Uscanga Aguilera}
\newcommand{\tarea}{Tarea 4}
\newcommand{\fecha}{13/3/2026}

\begin{document}
    \begin{center}
    {\large\textbf{UNIVERSIDAD DE GUANAJUATO}}\\[0.3cm]
    {\normalsize\textbf{DIVISIÓN DE CIENCIAS NATURALES Y EXACTAS}}\\
    {\normalsize\textbf{CAMPUS GUANAJUATO}}\\[1cm]

    {\Large\textbf{\tarea\ (\materia)}}\\[1cm]
    \end{center}

    \textbf{Nombre:} \alumno \hfill 
    \textbf{Fecha:} \fecha \hfill 
    \textbf{Calificación:} \rule{3cm}{0.4pt} \\[0.3cm]

    \begin{mdframed}[style=mdbluebox,frametitle={Ejercicio 1}]
        Un cuerpo, inicialmente en reposo $(\theta=0$ y $\omega=0$ cuando $t=0)$ es
        acelerado en una trayectoria circular de $1.3$m de radio, de acuerdo a la ecuacion
        $\alpha=120t^{2}-48t+16$. Encontrar la posición angular y la velocidad angular del
        cuerpo en función del tiempo, y las componentes tangencial y centrípeta de la
        aceleración.
    \end{mdframed}

    \begin{solucion}
        Velocidad angular. Integramos la aceleración angular
        \begin{align*}
            \omega(t)&=\int\alpha(t)dt\\
            &=\int(120t^{2}-48t+16)dt\\
            &=40^{3}-24t^{2}+16t+C
        \end{align*}
        Usando $\omega(0)=0$
        \begin{equation*}
            C=0
        \end{equation*}
        Por lo tanto
        \begin{equation*}
            \omega(t)=40t^{3}-24t^{2}+16t
        \end{equation*}
        Ahora posición angular. Integramos la velocidad angular
        \begin{align*}
            \theta(t)&=\int\omega(t)dt\\
            &=\int(40t^{3}-24t^{2}+16t)dt\\
            &=10t^{4}-8t^{3}+8t^{2}+C
        \end{align*}
        Usando $\theta(0)=0$
        \begin{equation*}
            C=0
        \end{equation*}
        Entonces
        \begin{equation*}
            \theta(t)=10t^{4}-8t^{3}+8t^{2}
        \end{equation*}
        La aceleracion tangencial es
        \begin{equation*}
            a_{t}=r\alpha
        \end{equation*}
        Sustituyendo
        \begin{equation*}
            a_{t}=1.3(120t^{2}-48t+16).
        \end{equation*}
        Así,
        \begin{equation*}
            a_{t}=156t^{2}-62.4t+20.8
        \end{equation*}
        Por ultimo, la aceleracion centripeta es
        \begin{equation*}
            a_{c}=r\omega^{2}
        \end{equation*}
        Sustituyendo
        \begin{equation*}
            a_{c}=1.3(40t^{3}-24t^{2}+16)^{2}
        \end{equation*}
    \end{solucion}

    \begin{mdframed}[style=mdbluebox,frametitle={Ejercicio 2}]
        Una pequeña masa $m$ se mueve sobre la mesa horizontal sin fricción. Está unida a
        una cuerda ligera que pasa por un pequeño agujero en la mesa. Del otro lado de la cuerda,
        inicialmente la partícula se mueve en una circunferencia de radio $r_{0}$ con velocidad
        tangencial $v_{0}$. Luego la cuerda se acorta y la partícula queda moviéndose en una
        circunferencia de radio $r$.
        \begin{enumerate}
            \item[(a)] Calcular el momento angular inicial de la partícula respecto al agujero.
            \item[(b)] Usar conservación del momento angular para encontrar la nueva velocidad tangencial $v(r)$
            \item[(c)] Determinar cómo cambia la energía cinética de la partícula.
            \item[(d)] Explicar por qué la energía cinética cambia aunque el momento angular se conserve.   
        \end{enumerate}
    \end{mdframed}

    \begin{solucion}
        \begin{enumerate}
            \item[(a)] Momengo angular inicial.\\
            El momento angular respecto al agujero es
            \begin{equation*}
                L=\vec{r}\times\vec{p}
            \end{equation*}
            En movimiento circular uniforme, su magnitud es
            \begin{equation*}
                L=mrv
            \end{equation*}
            Por lo tanto el momento angular inicial es
            \begin{equation*}
                L_{0}=mr_{0}v_{0}
            \end{equation*}
            \item[(b)] Velocidad tangencial usando conservacion del momento angular
            No hay torque externo respecto al agujero, por lo que
            \begin{equation*}
                L=\text{constante}
            \end{equation*}
            Entonces
            \begin{equation*}
                mr_{0}v_{0}=mrv(t)
            \end{equation*}
            cancelando $m$
            \begin{equation*}
                r_{0}v_{0}=rv(r)
            \end{equation*}
            Despejando $v(r)$
            \begin{equation*}
                v(r)=\frac{r_{0}}{r}v_{0}
            \end{equation*}
            Por lo tanto, la velocidad aumenta cuando el radio disminuye.
            \item[(c)] Cambio en la energía cinética
            La energia cinetica es
            \begin{equation*}
                K=\frac{1}{2}mv^{2}
            \end{equation*}
            Inicialmente
            \begin{equation*}
                K_{0}=\frac{1}{2}mv_{0}^{2}
            \end{equation*}
            Usando la velocidad obtenida
            \begin{equation*}
                K(r)=\frac{1}{2}m(\frac{r_{0}}{r}v_{0})^{2}\\
                K(r)=\frac{1}{2}mv_{0}^{2}(\frac{r_{0}}{r})^{2}
            \end{equation*}
            Como $r\lt r_{0}$, se tiene
            \begin{equation*}
                K(r)\lt K_{0}
            \end{equation*}
            Por lo tanto la energia cinetica aumenta.
        \end{enumerate}
    \end{solucion}
\end{document}